\subsection{Partial Inverses} 
Sometimes, a function is not one-to-one, but we can pretend. To do this, we restrict the \makeblank.
\begin{Example}
Let $f(x)=x^2$.
\begin{enumerate}\setabc
\item Give two pairs that show that $f$ is not one-to-one.\lbr\blankpair and \blankpair both have \makeblank[1in].
\item Let $g(x)=x^2,\ x\geq 0$. Is $g$ one-to-one? Why or why not?\makeblanknew
\item Given that $3^2=9$, what is $g\inv(9)$?\lbr$g\inv(9)=\makeblanksmall$
\item Let $h(x)=x^2,\ x\leq 0$. Is $h$ one-to-one? Why or why not?\makeblanknew
\item Given that $(-3)^2=9$, what is $h\inv(9)$?\lbr$h\inv(9)=\makeblanksmall$
\item Are $g$ and $h$ the same function? Why or why not?\makeblanknew
\end{enumerate}
\end{Example}
\noindent
This will be useful to us later in solving equations, and in creating new functions.